\documentclass{beamer}
% Computer Science 1 — shared Beamer preamble.
% \input this file after \documentclass{beamer} in every Monday lecture
% deck so all decks share one visual system. Do not fork per-week copies of
% this file — edit it once here and every deck picks up the change on next
% compile. Vendored from professional_minds/presentations/beamer/theme/
% preamble.tex (same pattern semester_kickoff_week's own vendored copy
% uses), kept local rather than \input across repos so this repo stays
% self-contained and usable without a professional_minds checkout.

\usetheme{Boadilla}
\usecolortheme{seahorse}
\setbeamertemplate{navigation symbols}{}
\setbeamertemplate{footline}[frame number]

% Speaker notes: \note{...} on any frame. Compiling with the default
% `pdflatex deck.tex` produces a slides-only PDF (notes are embedded but
% hidden). Compiling with
% `pdflatex "\PassOptionsToPackage{show}{pgfpages}\input{deck.tex}"`
% or uncommenting the line below produces an instructor copy with each
% slide's notes on a second page — see README.md for both commands.
% \setbeameroption{show notes on second screen=right}

\usepackage{pgfpages}

\newcommand{\SessionTitleSlide}[4]{%
  % #1 = Week/Day (e.g. "Week 2 -- Monday"), #2 = session title,
  % #3 = essential question, #4 = primary source citation
  \begin{frame}
    \titlepage
  \end{frame}
  \begin{frame}{Essential Question}
    \Large #3

    \vspace{1em}
    \normalsize\emph{#4}
  \end{frame}
}


\title{Week 10 Monday}
\subtitle{Lists: Collections That Grow}
\author{CS1}
\date{}

\begin{document}

\SessionTitleSlide{Week 10 -- Monday}{Lists: Collections That Grow}{How do we keep track of many things at once, and how do we critique the code that does it?}{Deitel \& Deitel, \emph{Intro to Python for Computer Science and Data Science}, Chapter 5, ``Sequences: Lists and Tuples''; AI Fluency Lens 10 (Critique) tie-in}

\begin{frame}[fragile]{The Problem: One Variable Per Thing}
\begin{verbatim}
colonist1 = "Ana"
colonist2 = "Bo"
colonist3 = "Cy"
print(colonist1)
print(colonist2)
print(colonist3)
\end{verbatim}
\begin{itemize}
    \item How would you print 50 names? Count them? Add one more later?
    \item We need \emph{one name} for \emph{many values}.
\end{itemize}
\note{Ask the class before showing the fix. Let them say ``a loop'' and then ``but the names are separate variables.'' Goal: feel the pain that a list removes.}
\end{frame}

\begin{frame}[fragile]{A List Holds Many Values in Order}
\begin{verbatim}
colonists = []                 # empty list
colonists.append("Ana")        # add one at a time
colonists.append("Bo")
colonists.append("Cy")
print(colonists)               # ['Ana', 'Bo', 'Cy']
print(len(colonists))          # 3
\end{verbatim}
\begin{itemize}
    \item \texttt{[]} makes an empty list; \texttt{.append(x)} adds to the end.
    \item \texttt{len(list)} is the number of items.
    \item \textbf{This is your gate:} build a list with \texttt{.append()} at least twice, then loop over it.
\end{itemize}
\note{Live code this. Students type along. Type the empty list, append three times, print after each append so they see it grow.}
\end{frame}

\begin{frame}[fragile]{Index: Position Starts at 0}
\begin{verbatim}
colonists = ["Ana", "Bo", "Cy"]
colonists[0]     # 'Ana'
colonists[2]     # 'Cy'
colonists[-1]    # 'Cy'  (last item)
colonists[3]     # IndexError: list index out of range
\end{verbatim}
\begin{itemize}
    \item Counting starts at 0; the last valid index is \texttt{len(list) - 1}.
    \item Negative indexes count from the end.
\end{itemize}
\note{Run the IndexError live; read the message aloud. Predict-then-run: what does colonists[1] give?}
\end{frame}

\begin{frame}[fragile]{Traverse: Do Something With Every Item}
\begin{verbatim}
for name in colonists:
    print("Welcome,", name)

total = 0
for name in colonists:
    total += 1
print("Headcount:", total)       # same as len(colonists)
\end{verbatim}
\begin{itemize}
    \item \texttt{for item in list:} visits each item in order.
    \item A traversal should \emph{report something real}: a roster, a count, a search result.
\end{itemize}
\note{This is the second half of the gate: build, then traverse and report. Connect to Week 5 loops.}
\end{frame}

\begin{frame}[fragile]{Changing a List}
\begin{verbatim}
crew = ["Ana", "Bo", "Cy"]
crew.insert(1, "Di")     # ['Ana', 'Di', 'Bo', 'Cy']
crew.remove("Bo")        # ['Ana', 'Di', 'Cy']
last = crew.pop()        # last == 'Cy'; crew == ['Ana', 'Di']
"Ana" in crew            # True
crew.index("Di")         # 1
\end{verbatim}
\begin{itemize}
    \item \texttt{.insert(i, x)}, \texttt{.remove(x)}, \texttt{.pop()}, \texttt{in}, \texttt{.index(x)}.
    \item Lists are \textbf{mutable}: these methods change the list itself.
\end{itemize}
\note{Predict-then-run: show each line, ask for the list afterward, then run. Mention that .remove raises ValueError if the item is missing.}
\end{frame}

\begin{frame}[fragile]{Sorting: Two Different Tools}
\begin{verbatim}
nums = [3, 1, 2]
sorted(nums)       # [1, 2, 3]   new list; nums unchanged
nums               # [3, 1, 2]
nums.sort()        # changes nums itself
nums               # [1, 2, 3]
x = nums.sort()    # x is None !
\end{verbatim}
\begin{itemize}
    \item \texttt{sorted(list)} returns a new list. \texttt{list.sort()} changes the list and returns \texttt{None}.
\end{itemize}
\note{This sets up mutation hazard 3. Ask: why would a method that changes a list return None?}
\end{frame}

\begin{frame}{Critique Mode: Find the Weakness}
\begin{itemize}
    \item Monday Moments this week: \textbf{Lens 10, Critique} --- find weaknesses in work before trusting it.
    \item Reviewing code is the same habit: read it, predict what it does, then test.
    \item Next: three snippets that \emph{look} fine and are not.
\end{itemize}
\note{Pairs for the next three slides. Give two minutes per snippet: predict the output, then run it, then name the weakness in one sentence.}
\end{frame}

\begin{frame}[fragile]{Hazard 1: A Second Name Is Not a Copy}
\begin{verbatim}
crew = ["Ana", "Bo", "Cy"]
backup = crew
backup.append("Di")
print(crew)        # what prints?
\end{verbatim}
\pause
\begin{verbatim}
# ['Ana', 'Bo', 'Cy', 'Di']  -- crew changed too!
backup = crew.copy()      # or list(crew)
\end{verbatim}
\begin{itemize}
    \item \texttt{backup = crew} makes two names for \emph{one} list.
    \item \texttt{.copy()} makes a new list.
\end{itemize}
\note{Draw two name tags on one box, then two boxes. If time, show id(crew) == id(backup).}
\end{frame}

\begin{frame}[fragile]{Hazard 2: Don't Remove While You Loop}
\begin{verbatim}
crew = ["Ana", "Bo", "Bo", "Cy"]
for name in crew:
    if name == "Bo":
        crew.remove(name)
print(crew)        # what prints?
\end{verbatim}
\pause
\begin{verbatim}
# ['Ana', 'Bo', 'Cy']  -- one Bo survived!
crew = [n for n in crew if n != "Bo"]   # build a new list
\end{verbatim}
\begin{itemize}
    \item Removing shifts the items, so the loop skips one.
    \item Safer: loop over a copy, or build a new list.
\end{itemize}
\note{List comprehension is a preview; the safe version for Week 10 is for name in crew.copy() or build a second list with append.}
\end{frame}

\begin{frame}[fragile]{Hazard 3: sort() Returns None}
\begin{verbatim}
scores = [90, 72, 85]
ranked = scores.sort()
print(ranked)       # what prints?
\end{verbatim}
\pause
\begin{verbatim}
# None
ranked = sorted(scores)     # or: scores.sort(); then use scores
\end{verbatim}
\begin{itemize}
    \item Methods that change a list in place usually return \texttt{None}.
\end{itemize}
\note{Ties back to the sorting slide. Ask pairs: what would a reviewer write in a code comment here?}
\end{frame}

\begin{frame}{Your Gate: The Census}
\begin{itemize}
    \item Your world's population, inventory, case load, or customers becomes a list.
    \item \textbf{Quick Check (pass/fail):}
    \begin{itemize}
        \item The list holds several distinct entities from your world.
        \item It is built with \texttt{.append()} at least twice (not one typed-out literal).
        \item At least one loop or \texttt{len()} reports something real.
    \end{itemize}
    \item World Bible: one line --- what you passed, what broke.
\end{itemize}
\note{Frontier Settlement: colonists. Investigation Bureau: suspects. Starship Log: crew. Small Business: customers.}
\end{frame}

\begin{frame}[fragile]{Work Time: Start Your Census}
\begin{verbatim}
colonists = []
colonists.append("...")     # at least two appends
colonists.append("...")
for c in colonists:
    print(c)
print("Headcount:", len(colonists))
\end{verbatim}
\begin{itemize}
    \item Work for about 10 minutes; ask a neighbor to critique before you ask me.
\end{itemize}
\note{Circulate. Common issues: forgetting the empty list before append, calling append on a string, misspelling the variable.}
\end{frame}

\begin{frame}{Wrap-Up and Preview}
\begin{itemize}
    \item Lists: \texttt{[]}, \texttt{.append}, \texttt{len}, index, \texttt{for}, \texttt{in}.
    \item Hazards: aliasing, removing while looping, \texttt{.sort()} returns \texttt{None}.
    \item \textbf{Week 11:} dictionaries --- look up an item by name or ID instead of searching the whole list.
    \item Gate and optional bonus practice are due \textbf{Friday 11:59 PM}.
\end{itemize}
\note{Ask one student to show their list. Remind them about pair programming on the gate.}
\end{frame}

\end{document}
